\documentclass[preprint,12pt]{elsarticle}
\usepackage{amsmath,amssymb,amsfonts,amsthm}
\usepackage{xcolor}
\usepackage{graphicx}
\usepackage{booktabs}
\usepackage{bm}
\usepackage{url}
\DeclareUrlCommand\path{\urlstyle{rm}}
\providecommand{\href}[2]{#2}

\newcommand{\tikzdir}{figures/tikz}
% Shared PGFPlots/TikZ style for the figures of this project.
% Modelled on the figures in M. Ambarkutuk and P. E. Plassmann, "Modeling and analysis of dispersive propagation of structural waves for
% vibro-localization", Sensors 24(23):7744, 2024 (read in full): boxed axes with outward ticks, DejaVu Sans text, thin lines, two-colour
% proposed-vs-baseline convention (red = proposed, black = baseline; the 2023 Sensors paper uses blue/brown for its PDF/CDF pair),
% axis labels written "Quantity [unit]", frame-boxed legends inside the axes, viridis maps, red x / black + markers, (a)(b)(c) panel labels below.
\usepackage{xcolor}
\usepackage{tikz}
\usepackage{pgfplots}
\usepgfplotslibrary{groupplots}
\usetikzlibrary{calc,patterns,positioning}
\usepackage{xstring}
\pgfplotsset{compat=1.16}
\usepackage[scaled=0.95]{DejaVuSans}
\usepackage{sansmath}
\newcommand{\figfont}{\fontencoding{T1}\fontfamily{DejaVuSans-TLF}\selectfont\sansmath}
\providecommand{\tikzdir}{.}

\definecolor{ambRed}{RGB}{255,0,0}
\definecolor{ambBlack}{RGB}{0,0,0}
\definecolor{ambBlue}{RGB}{31,119,180}
\definecolor{ambBrown}{RGB}{140,86,75}
\definecolor{ambGreen}{RGB}{44,160,44}
\definecolor{ambGrey}{RGB}{217,217,217}
\definecolor{ambPass}{RGB}{207,232,213}
\definecolor{ambFail}{RGB}{243,198,198}

% diverging maps (ColorBrewer RdBu reversed, PiYG) and viridis/magma are built in or defined here
\pgfplotsset{
  colormap={rdbu}{rgb255=(5,48,97) rgb255=(33,102,172) rgb255=(67,147,195) rgb255=(146,197,222) rgb255=(209,229,240) rgb255=(247,247,247)
                  rgb255=(253,219,199) rgb255=(244,165,130) rgb255=(214,96,77) rgb255=(178,24,43) rgb255=(103,0,31)},
  colormap={piyg}{rgb255=(142,1,82) rgb255=(197,27,125) rgb255=(222,119,174) rgb255=(241,182,218) rgb255=(253,224,239) rgb255=(247,247,247)
                  rgb255=(230,245,208) rgb255=(184,225,134) rgb255=(127,188,65) rgb255=(77,146,33) rgb255=(39,100,25)},
}

\pgfplotsset{
  colormap={viridisa}{rgb255=(68,1,84) rgb255=(72,40,120) rgb255=(62,74,137) rgb255=(49,104,142) rgb255=(38,130,142) rgb255=(31,158,137)
                      rgb255=(53,183,121) rgb255=(109,205,89) rgb255=(180,222,44) rgb255=(253,231,37)},
  colormap={magmaa}{rgb255=(0,0,4) rgb255=(28,16,68) rgb255=(79,18,123) rgb255=(129,37,129) rgb255=(181,54,122) rgb255=(229,80,100)
                    rgb255=(251,135,97) rgb255=(254,194,135) rgb255=(252,253,191)},
}
\pgfplotsset{
  amb/.style={
    scale only axis, axis line style={line width=0.5pt}, tick align=outside, tick pos=left,
    major tick length=2.5pt, tick style={line width=0.5pt},
    label style={font=\figfont\footnotesize}, tick label style={font=\figfont\scriptsize}, title style={font=\figfont\footnotesize},
    legend style={font=\figfont\tiny, draw=black!25, fill=white, fill opacity=0.85, text opacity=1, rounded corners=1pt,
                  inner sep=2pt, nodes={inner sep=1pt}, legend cell align=left},
    every axis plot/.append style={line width=0.8pt, mark size=1.3pt},
    ymajorgrids=false, xmajorgrids=false,
  },
  % method styles shared by the toy figures
  known/.style={ambBlack, dashed, mark=square*, mark options={solid}},
  oracle/.style={ambRed, mark=*},
  kernel/.style={ambBlue, mark=triangle*, mark options={solid}},
  kernelnr/.style={ambBrown, mark=triangle*, mark options={rotate=180,solid}},
  shrunk/.style={ambGreen, mark=diamond*, mark options={solid}},
}
\pgfplotsset{barleg/.style={legend image code/.code={\draw[#1] (0cm,-0.08cm) rectangle (0.28cm,0.08cm);}}}
\newcommand{\subcap}[2]{\node[font=\figfont\footnotesize, anchor=north] at (#1) {#2};}
   % PGFPlots/TikZ style of the figures (DejaVu Sans text, boxed axes, red/black/blue/brown series)
\usepackage{adjustbox}
\newlength{\figw}
% ---- shared macros -------------------------------------------------------
\newcommand{\TODO}[1]{\textcolor{red}{\textbf{[TODO: #1]}}}
\newcommand{\todocell}{\textcolor{red}{\textbf{TODO}}}

% placeholder figure: a framed box that is clearly marked as not yet produced
\newcommand{\todofig}[2]{%
  \begin{figure}[!t]
    \centering
    \fbox{\parbox[c][2.6cm][c]{0.92\columnwidth}{\centering
      \textcolor{red}{\textbf{TODO FIGURE}}}}
    \caption{#2\ \textcolor{red}{\textbf{[PLACEHOLDER --- not yet produced]}}}
    \label{fig:#1}
  \end{figure}}

\newcommand{\runin}[1]{\par\smallskip\noindent\textit{#1}\ }
\DeclareMathOperator*{\argmax}{arg\,max}
\DeclareMathOperator*{\argmin}{arg\,min}
\newcommand{\R}{\mathbb{R}}
\newcommand{\ray}{\mathbf{d}}
\newcommand{\nrm}{\mathbf{n}}
\newcommand{\vars}{\sigma}
\newcommand{\Sgrid}{\mathcal{S}}
\newtheorem{proposition}{Proposition}
\newtheorem{remark}{Remark}


\journal{Preprint}

\begin{document}
\begin{frontmatter}

\title{Calibrated Multi-View-Consistent Depth Sampling with a Diffusion Model}

\author{Murat Ambarkutuk}
\ead{murat@ambarkutuk.com}
\affiliation{organization={Independent Researcher},
            city={Shelton},
            state={CT},
            country={USA}}

\begin{abstract}
Single-image depth networks leave global scale free, fail regionally, and return a point estimate against which no covariance can be tested.
We study a generative alternative: a latent diffusion model (SDXL with low-rank adaptation) fine-tuned to sample an affine-invariant per-view depth from the RGB image alone.
Its unguided samples are the object of a leave-view-out fault test: a Gaussian-splat (2DGS) surface is built from the other posed views only, never view $j$,
a per-view scale and shift is fitted between the samples and the render, and the residual is tested against the sample covariance plus the render covariance with an empirically calibrated threshold.
Guided sampling, in which the likelihood gradient toward the leave-view-out render acts in the last denoising steps, is intended as the multi-view-consistent product; training is planned to be semi-self-supervised, with ground-truth depth where it exists and render-and-compare elsewhere.
So far we have (i) a synthetic check of the statistic, in which fitting the scale and shift costs almost no power when the covariance is known (minimum detectable fault $0.88\sigma$ against $0.89\sigma$) and an estimated covariance accounts for most of the power loss,
and (ii) a supervised first stage of the model, trained for 2{,}000 of a planned 10{,}000 steps on BlendedMVS and Hypersim, whose validation loss fell monotonically; on three held-out views of one ETH3D scene the affine-fitted AbsRel (fitted to sparse laser depth) was between $0.10$ and $0.11$ depending on checkpoint and number of sampler steps, whereas on one of those views, the only one where we scored it, the untrained base gave $0.38$ (trained: $0.11$--$0.14$).
That check is a sanity test, not a result.
The semi-self-supervised stage, guided sampling, and the calibration and fault-injection experiments on real model samples have not been run, and no claim is made about them.
\end{abstract}

\begin{keyword}
diffusion models \sep monocular depth estimation \sep uncertainty calibration \sep Gaussian splatting \sep multi-view consistency \sep fault detection
\end{keyword}

\end{frontmatter}

\section{Introduction}
\label{sec:intro}
\TODO{Motivation: monocular depth is useful for multi-view reconstruction, but a single image leaves global scale free and
the network fails regionally; a point estimate has no covariance to test against.}

\runin{Contribution (planned).}
(i) A diffusion model that samples an affine-invariant depth posterior for a view from its RGB image alone (unguided mode).
(ii) Sampling-time guidance from a Gaussian-splat surface fitted to the \emph{other} views only (guided mode), whose output is the
multi-view-consistent depth product.
(iii) Semi-self-supervised training: ground-truth depth where it exists, render-and-compare elsewhere, with the labelled fraction varied.
(iv) The spread of unguided samples as a covariance for a leave-view-out fault test; calibration is reported, not only accuracy.
(v) A fault-injection protocol for a learned prior that may explain a real fault as plausible depth.
\TODO{Reduce to the contributions the experiments support.}

\section{Related Work}
\label{sec:related}
\runin{Diffusion monocular depth.}
Marigold~\cite{ke2024marigold} fine-tunes a latent image diffusion model on synthetic data to predict affine-invariant depth and
ensembles members after aligning their scale and shift; its ensemble spread is not evaluated against error.
\runin{Multi-view diffusion for depth.}
MVGD~\cite{guizilini2025mvgd} generates novel-view images and depth from posed views with a pixel-space diffusion model trained
from scratch; it outputs no uncertainty, and for stereo and video depth the target image is also a conditioning input.
\runin{Guidance for consistency.}
DepthSync~\cite{depthsync2025} guides a video diffusion depth model at inference with scale and geometry losses between neighbouring frames;
the target frame's own prediction enters the guide, and it reports no uncertainty.
\runin{Depth rectification in splatting.}
AffineGS~\cite{tse2026affinegs} corrects a depth network's output toward the splat render of the same view with a per-pixel affine field;
it has no uncertainty and no leave-view-out step.
\TODO{Fuller literature pass (video depth diffusion, uncertainty for depth, fault detection and exclusion, splatting with depth priors).
Claims in this section rest on the papers as read; see \path{refs/EVIDENCE.md}.}

\section{Method}
\label{sec:method}

\subsection{Problem setting}
\label{sec:method:problem}
There are $K$ depth sensors. Sensor $s$ produces depth images
$D_s^t:\Omega_s\to\R_{>0}$ at times $t$, with intrinsics $\mathbf{K}_s$ and a
believed pose $\mathbf{T}^{t}_{s}$. A subset $\mathcal{F}\subset\{1,\dots,K\}$ of
sensors is \emph{faulty}: its depth, intrinsics or pose deviate persistently
from the truth in a way that is not zero-mean noise (extrinsic drift, depth
bias, multipath; Section~\ref{sec:setup:faults}). Neither $\mathcal{F}$ nor
$|\mathcal{F}|$ is known; we assume $|\mathcal{F}|\le F$ with $F<K/2$ and
return to this assumption in Section~\ref{sec:discussion}. The goal is a fused
surface that is close to the reconstruction from the healthy sensors alone,
together with the identity of the excluded sensors.

\subsection{Shared surface and per-ray noise model}
\label{sec:method:noise}
The shared surface is a set of $N$ Gaussians with means $\boldsymbol\mu_i$,
covariances $\boldsymbol\Sigma_i$ and opacities $\alpha_i$. $N$ is fixed and
optimised with the relocation scheme of 3DGS-MCMC \cite{gs_mcmc}, so densification
cannot create a second layer that absorbs a faulty sensor's surface. Because
photometric supervision is absent, we add an accumulated-alpha loss on valid and
invalid pixels and a cap on splat anisotropy; planar splats follow
\cite{gs_2dgs}. \TODO{Specify loss weights after the dev-scene sweep.}

For a ray of length $r$ that meets a locally planar surface at incidence angle
$\theta$ (angle between ray and surface normal), a range noise $\sigma_r$ and a
pointing noise $\sigma_\varphi$ give a depth \emph{residual} variance
\begin{equation}
  v_{\mathrm{res}}(r,\theta)=\sigma_r^{2}+\bigl(r\,\sigma_\varphi\tan\theta\bigr)^{2}.
  \label{eq:vres}
\end{equation}
This follows from $\mathrm{d}r/\mathrm{d}\theta=r\tan\theta$ for a plane at
perpendicular distance $h$, where $r=h/\cos\theta$. It is the variance of the
discrepancy between the measured range and the intersection of the nominal ray
with the surface; it is the right weight for a depth residual along the ray. It is
not the covariance of the 3D point: in the ray frame $(\ray,\mathbf{u},\mathbf{v})$
the point covariance is $\mathrm{diag}(\sigma_r^2,(r\sigma_\varphi)^2,(r\sigma_\varphi)^2)$,
and its component along the surface normal is
\begin{equation}
  v_{n}(r,\theta)=\sigma_r^{2}\cos^{2}\theta+r^{2}\sigma_\varphi^{2}\sin^{2}\theta,
  \label{eq:vn}
\end{equation}
which stays bounded as $\theta\to\pi/2$. We therefore use \eqref{eq:vres} to
weight residuals and \eqref{eq:vn} (with the tangential scale stretched by
$1/\cos\theta$ in the plane of incidence) to initialise and floor splat
covariances. Equation~\eqref{eq:vres} diverges at grazing incidence and its
linearisation fails once $\sigma_\varphi\tan\theta\gtrsim 0.3$, so we use
the regularised form
\begin{equation}
  \tilde v_{\mathrm{res}}=\operatorname{clip}\!\Bigl(\sigma_{r,\mathrm{eff}}^{2}
  +r^{2}\sigma_{\varphi,\mathrm{eff}}^{2}\,\tfrac{\sin^{2}\theta}{\cos^{2}\theta+\epsilon^{2}},\;
  v_{\min},\,v_{\max}\Bigr),
  \label{eq:vreg}
\end{equation}
with $\epsilon\approx\cos 80^{\circ}$, a hard gate at $\theta>85^{\circ}$, and a
beam-footprint term added to $\sigma_\varphi^2$ for LiDAR and ToF. For stereo
and structured-light sensors the range noise is $\sigma_z=z^{2}\sigma_d/(fb)$
\cite{noise_khoshelham,noise_matthies}; ToF noise grows as signal amplitude falls
\cite{noise_lange,noise_sarbolandi}. ToF multipath and flying pixels are biases
rather than variance and are masked by amplitude and depth-gradient tests
\cite{noise_freedman}. We use the per-ray variance in three places: as a
likelihood weight $\tfrac12 r^2/\tilde v_{\mathrm{res}}$, to initialise splat
scales, and as a soft floor on splat scale. It is not used as a fixed
covariance prior, because the posterior variance shrinks as views accumulate.

\subsection{Consensus rendering and visibility-aware residuals}
\label{sec:method:consensus}
Let $\hat D_{-s}^{t}$ be the depth rendered into sensor $s$'s view from a surface
fitted \emph{without} sensor $s$ (a warm-started refit of about $200$
iterations), with accumulated alpha $A_{-s}^{t}$ and variance $v_{-s}$. Fitting
without $s$ prevents self-influence: a joint fit can absorb a sensor's own bias
and make its residuals look small. A pixel $u$ enters the residual only if it is
valid, $A_{-s}^{t}(u)>\tau_A$, and the rendered first return agrees with the
sensor's first return within a visibility gate, so that occluded regions do not
count as disagreement. The normalised residual and per-sensor statistic are
\begin{equation}
  \rho_s(u)=\frac{D_s(u)-\hat D_{-s}(u)}{\sqrt{\tilde v_{\mathrm{res},s}(u)+v_{-s}(u)}},
  \qquad
  T_s=\frac{1}{|\Omega_s^{\mathrm{vis}}|}\sum_{u\in\Omega_s^{\mathrm{vis}}}\rho_s(u)^{2}.
  \label{eq:Ts}
\end{equation}
Pixel residuals are spatially correlated, so we do not use the nominal
$\chi^2$ quantile as a threshold; the threshold $\tau$ is calibrated on
fault-free dev runs to a 5\% per-run false-elimination rate
(Section~\ref{sec:setup:protocol}). We also record Moran's~I of $\rho_s$ as a
diagnostic: a coherent fault leaves spatially structured residuals, whereas
sparse outliers do not. \TODO{Report Moran's~I figure (Fig.~\ref{fig:residual-structure}).}

\subsection{Corrected admission criterion}
\label{sec:method:admission}
Ambarkutuk \emph{et~al.}\ \cite{flt_ambarkutuk} give each sensor $i$ a PDF $f_i$
over a common grid $\Sgrid$, form the product over a consensus set $C$, and grow $C$
greedily from the best pair using an entropy test. In our reading of the
published text, the prose and Eq.~(14) seed the set with the pair of
\emph{maximum} entropy after fusion, whereas the corresponding line of
Algorithm~1 tests $h_{ij}\le h^{*}$ with $h^{*}$ initialised to $-\infty$, so the
seed branch never executes; the expansion test accepts a sensor when entropy
does not decrease, which is the opposite of the stated aim of sharpening the
joint PDF, and it depends on sensor order.
\TODO{Re-check the direction of the inequalities against the typeset equations in the
  published PDF; our reading used a text rendering in which math was flattened.}

A more basic problem is that entropy of a normalised product cannot detect a
biased sensor. For one-dimensional Gaussians $\mathcal{N}(m_C,\sigma_C^2)$ and
$\mathcal{N}(m_k,\sigma_k^2)$, the normalised product has variance
$\sigma_C^2\sigma_k^2/(\sigma_C^2+\sigma_k^2)$ regardless of the offset
$m_C-m_k$, so
\begin{equation}
  H(C\cup k)-H(C)=-\tfrac12\log\!\bigl(1+\sigma_C^{2}/\sigma_k^{2}\bigr)<0
  \label{eq:dH}
\end{equation}
for every sensor, healthy or faulty; a sharp biased sensor lowers entropy the
most. Entropy measures the precision of a consensus, not its consistency.

We therefore separate the two roles. Let $F_C=\prod_{i\in C}f_i$,
$Z_C=\int_{\Sgrid}F_C$, $p_C=F_C/Z_C$ and $H(C)=-\!\int p_C\log p_C$, all evaluated in
log space on the grid. Sensor $k$ is \emph{compatible} with $C$ when its Bayes
factor against a uniform null on $\Sgrid$ is at least one:
\begin{equation}
  \Lambda_k(C)=|\Sgrid|\int_{\Sgrid}p_C f_k\,\mathrm{d}x=|\Sgrid|\,\frac{Z_{C\cup k}}{Z_C}\;\ge\;1.
  \label{eq:lambda}
\end{equation}
For Gaussians, $\Lambda_k(C)\ge 1$ is equivalent to
$\tfrac12\nu^{2}/(\sigma_C^2+\sigma_k^2)\le \log\!\bigl(|\Sgrid|/\sqrt{2\pi(\sigma_C^2+\sigma_k^2)}\bigr)$
with $\nu=m_C-m_k$, so the offset enters exponentially.
Algorithm~\ref{alg:admission} seeds with the pair maximising $Z_{ij}$ (and
requires it to pass the test), expands with the compatible sensor of largest
$\Lambda_k$, which removes the order dependence, and finally returns the
minimum-entropy set among admissible sets with at most $F$ removals. The grid
volume $|\Sgrid|$ is part of the test and must be fixed and reported. For $K<3$ or
when no admissible set satisfies the removal budget, the algorithm returns
\emph{no consensus} and excludes nobody.

\begin{algorithm}[t]
\caption{Admission-based consensus (corrected criterion)}
\label{alg:admission}
\begin{algorithmic}[1]
\REQUIRE PDFs $\{f_i\}_{i=1}^{K}$ on grid $\Sgrid$, removal budget $F<K/2$
\STATE $(i^\star,j^\star)\gets\argmax_{i\ne j}Z_{\{i,j\}}$
\IF{$\Lambda_{j^\star}(\{i^\star\})<1$}
  \RETURN no consensus
\ENDIF
\STATE $C\gets\{i^\star,j^\star\}$
\WHILE{$\exists\,k\notin C:\ \Lambda_k(C)\ge 1$}
  \STATE $k^\star\gets\argmax_{k\notin C}\Lambda_k(C)$; \ $C\gets C\cup\{k^\star\}$
\ENDWHILE
\STATE Among admissible sets $C'$ with $K-|C'|\le F$, \textbf{return} $\argmin_{C'}H(C')$
\end{algorithmic}
\end{algorithm}

\begin{proposition}[Relation to leave-one-out and solution-separation tests]
\label{prop:loo}
Let sensor $k$ measure a state $x$ directly with noise covariance $R_k$, let
$\hat x_{-k}$ be the estimate without it with covariance $P_{-k}$, and let
$\nu=z_k-\hat x_{-k}$ with $S=R_k+P_{-k}$. Then the solution-separation statistic
$d^{\top}\mathrm{Cov}(d)^{-1}d$, with $d=\hat x_{\mathrm{all}}-\hat x_{-k}$,
equals the leave-one-out innovation statistic $\nu^{\top}S^{-1}\nu$. For Gaussian
PDFs the negative log-evidence of $k$ under the consensus,
$-\log\!\int p_C f_k=\tfrac12\nu^{\top}S^{-1}\nu+\tfrac12\log\det(2\pi S)$, is the leave-one-out statistic plus a covariance-volume term.
\end{proposition}
\begin{proof}[Sketch]
$d=P_{-k}S^{-1}\nu$ and $\mathrm{Cov}(d)=P_{-k}-P_{\mathrm{all}}=P_{-k}S^{-1}P_{-k}$, so
\begin{align*}
d^{\top}\mathrm{Cov}(d)^{-1}d
 &=\nu^{\top}S^{-1}P_{-k}\bigl(P_{-k}S^{-1}P_{-k}\bigr)^{-1}P_{-k}S^{-1}\nu\\
 &=\nu^{\top}S^{-1}\nu .
\end{align*}
The evidence is the Gaussian density of $\nu$ under $\mathcal{N}(0,S)$.
\end{proof}
\begin{remark}
Under equal covariances and Gaussian residuals, ranking sensors by this evidence is
therefore the same as ranking by the leave-one-out chi-square statistic, and
the entropy-based criterion can add value only where noise is heteroscedastic,
residuals are non-Gaussian or multimodal, or occlusion and grazing geometry
break the chi-square assumptions. We test exactly this in H2
(Section~\ref{sec:setup:hyp}) and treat a tie as a reportable outcome.
\end{remark}

\subsection{Elimination, refit and persistence}
\label{sec:method:elim}
We apply the admission test to the per-sensor depth posteriors induced by the
consensus renders. Procedure: (i)~fit jointly; (ii)~compute $T_s$ and the
admission statistics by warm-started leave-one-sensor-out refits; (iii)~exclude
the sensor with the strongest evidence of inconsistency if it exceeds the
calibrated threshold; (iv)~refit and repeat until no sensor is flagged or
fewer than three remain. To separate sensor faults from dynamics we track the
per-sensor, per-frame residual energy $e_{s,t}=\overline{\rho_s^2}$: a
persistent rise in one sensor while others stay flat is flagged by a CUSUM
test as a sensor fault, whereas a localised rise visible to several sensors in
few frames is treated as a dynamic object and masked.
\TODO{Dynamics handling is optional and first to be cut if the schedule slips;
  state its status in the final version.}

\section{Experimental Design and Hypotheses}
\label{sec:setup}
The experiments ask when the gate helps and when it does not, and they start by measuring the
depth model. All use public datasets, public pretrained models and public code; the protocol,
hypotheses and thresholds are fixed and tagged in the code repository before any test scene is run.
No experiment has been run for this version.

\subsection{Design in Brief}
\label{sec:setup:expsetup}
\emph{Data.} ETH3D \cite{data_eth3d}, with laser-scan depth and calibrated poses, is primary and is
where the residuals are measured; ScanNet++ \cite{data_scannetpp}, subject to its access terms, is
secondary and contains many reflective and transparent surfaces. Poses come from the datasets.
\emph{Depth model.} Depth Pro \cite{est_depthpro}, one prediction per image; Metric3D v2 and
UniDepthV2, which output a confidence, supply the confidence baseline. \emph{Reconstruction.}
Depth-supervised splatting of 2DGS surfels with gsplat, with the appearance block of
Section~\ref{sec:method:setting}, run with and without the gate. \emph{Faults.} Natural failures
labelled against the laser depth on connected regions, classed as reflective or glass, textureless
wall, or other; and injected regional faults (offset, plane replacement, blurred copy and mirrored
depth on 2--40\% of a view, at 2--16$\sigma$), placed on appearance masks or random superpixels,
either independent per view or repeated across views. \emph{Baselines.} SpotLessSplats
\cite{gs_spotless}, the RobustNeRF loss \cite{gs_robustnerf} and robust depth losses (Huber,
Geman--McClure) in the fit; as view tests, the per-view aligned RMS residual against the other
views' render and the depth model's confidence; no gate and an oracle gate. \emph{Metrics.}
Detection power at 5\% per-view false alarm, false-alarm and untestable rates, held-out laser-depth
AbsRel, and surface accuracy, completeness and F-score. Details are in Supplement~\ref{sup:setup}.

\subsection{Hypotheses}
\label{sec:setup:hyp}
Each hypothesis states what would falsify it, and several state where the gate is not expected to
help. Thresholds match the pre-registered configuration files.
\begin{description}
  \item[H1] \emph{(residuals first; go/no-go).} Fit $\sigma(\cdot)$, $\ell$, $\nu$ and tails of the
        aligned Depth Pro residuals on ETH3D development scenes; characterise natural regional
        failures. \emph{Go} if $\ge10\%$ of views carry a failure covering $\ge2\%$ of pixels at
        $>3\sigma$ and keeping those views lowers the fused-surface F-score by $\ge2$ points.
        \emph{No-go} (pivot to an analysis paper) below 5\% of views or below 1 point.
  \item[H2] \emph{(repetition).} Cross-view correlation of aligned errors at the same laser surface
        point: $\ge0.5$ on textureless regions, $\le0.3$ on reflective, $\le0.2$ on textured.
        \emph{Falsified} if the textureless value is below 0.5 or not at least 0.2 above the reflective one.
  \item[H3] \emph{(detection).} At 5\% false alarm, on reflective-region failures, the gate beats the
        best cheap or robust baseline by $\ge15$ points of power. \emph{Falsified} if the gain is
        below 15 points or its interval includes zero. No gain is predicted on textureless failures.
  \item[H4] \emph{(reconstruction).} With the gate, AbsRel and F-score improve by $\ge10\%$ relative
        over no gate and the best robust-splatting baseline on failure scenes of ETH3D and ScanNet++,
        and change by $\le2\%$ on clean scenes. \emph{Falsified} if the paired interval includes zero
        on failure scenes, or clean scenes degrade by more than 2\%.
  \item[H5] \emph{(calibration).} Raw mean $T_j/(n-2)>1.2$ on clean views; after deflation or
        threshold calibration, the test-scene false-alarm rate is within $[0.5\alpha,1.5\alpha]$.
        \emph{Falsified} if it leaves that band.
  \item[H6] \emph{(ablations).} Joint surface $-10$ points; keeping tested-view-only surfels $-20$
        points on large region failures; unwhitened $-5$ points; rendered blocks at least as calibrated
        and powerful as the FisherRF-style and Bayes' Rays-style uncertainties. \emph{Falsified} if any
        ablation is within its margin.
\end{description}
H1 is a go/no-go and runs first; H2 predicts a blind spot rather than a success; H3 and H4 are the
claims that would justify the method; H5 tests the stated approximation; H6 tests the design choices.

\paragraph{Breakdown scenarios} Reported as honest negatives: errors repeated in all views
(textureless walls, H2); low overlap, where pixels are untestable; faulty views that supply most of a
region; a misspecified measurement model; and faults or misfits beyond the linearised fusion.

\section{Results}
\label{sec:results}
\textcolor{red}{\textbf{No experiment has been run for this version. Every table and figure is a
placeholder and must be regenerated by the experiment code before submission.}} The planned
tables and figures, one per hypothesis, are in Supplement~\ref{sup:results}; the main-text results
will report the regime map (H1), the consequence for trajectory and surface (H2) and the price of
monitoring (H6).

\section{Conclusions and Future Work}
\label{sec:conclusion}
We have proposed a view-level gate that decides which depth sources to trust before pose
graph optimisation and fusion. Depth sources are the outputs of an ensemble of monocular
estimators on RGB images; the gate compares them after a leave-one-source-out alignment,
on pixels visible in both views, and admits a source only if its mean log Bayes factor
against the consensus passes a calibrated threshold. The admission test corrects the
entropy criterion of our earlier Byzantine sensor elimination, which could not
distinguish an agreeing source from a sharp but biased one, and returns \emph{no
consensus} when attribution is not identifiable. A test bed, consisting of a monocular
depth ensemble, a partly end-to-end pose graph and fusion, provides the depth maps and
the downstream task; MASt3R-SfM and VGGT are reference points, and robust kernels, the
COLMAP geometric consistency filter and plain voting are the baselines the gate must
beat. The empirical case rests on hypotheses H1--H5, none of which has been tested in
this version. \TODO{Rewrite the conclusion around the measured results; if H2 fails,
reframe around the gate with the leave-one-out chi-square statistic and keep the entropy
relation as an equivalence result.}

\TODO{Interpretation of the results, to write after the experiments: for each of H1--H5
state whether it was supported, partly supported or falsified, and why; state from H2
whether the evidence criterion adds anything beyond the leave-one-out chi-square test
(Proposition~\ref{prop:loo}); and from H1 in which overlap and clustering regimes the gate
beats the pixel filters and robust kernels and where it ties.}

\subsection{Limitations and Failure Conditions}
\label{sec:disc:limits}
The gate is built on a consensus and fails where the consensus is wrong or not
identifiable. First, it assumes that fewer than half of the sources of any view are
faulty. Monocular estimators trained on overlapping data tend to fail on the same
content, so a majority can fail together; the gate is then expected to be no better
than fusing everything, and it can exclude the healthy minority. Second, with $M=2$
the residuals are symmetric and attribution is not identifiable; only detection is
possible. Third, the gate targets structure that the scale-and-shift alignment cannot
remove; a pure scale or shift error is absorbed by the alignment, in the gate and in the
pose graph, and is not detected. Fourth, the gate uses provisional poses to warp depth
between views, so poor provisional poses can corrupt it. Fifth, faults below the
per-pixel noise can be detected only by aggregating over many pixels, so there is a
detection floor that depends on overlap and pixel count. Sixth, excluding a source
lowers completeness wherever only that source covered a region. Seventh, the injected
faults are controlled surrogates of natural failures; natural failures are scored against
a laser reference on a small number of scenes, and the scenes are static. Finally,
leave-one-source-out refits and the estimator ensemble cost memory and time; the test bed
is an offline pipeline.
\TODO{Update with observed failures once experiments run.}

\subsection{Future Work}
\label{sec:disc:future}
Planned extensions, none yet started, are: training the gate through the differentiable
back-end; running it online on video under a per-frame latency budget; handling dynamic
scenes; and applying the same gate to depth from physical sensors (LiDAR, time-of-flight,
active stereo), where the sources are sensors rather than estimators and the fault
classes include extrinsic drift and multipath. \TODO{Confirm which of these remain in
scope.}

\label{endmaintext}

\section*{Declarations}
\textbf{Funding.} This research received no external funding.
\textbf{Data and code availability.} \TODO{State what is released. Restricted datasets are not redistributed.}
\textbf{Use of AI tools.} \TODO{Author to state whether and how generative AI tools assisted the writing.}

\bibliographystyle{elsarticle-num}
\bibliography{refs/references}

\end{document}
